\documentclass[12pt,a4paper]{article}
\usepackage[UTF8]{ctex}
\usepackage{amsmath,amssymb}
\usepackage{geometry}
\usepackage{listings}
\usepackage{xcolor}
\usepackage{tcolorbox}
\usepackage{enumitem}
\usepackage{hyperref}

\geometry{margin=2.5cm}

% C++ 代码高亮设置
\lstset{
	language=C++,
	basicstyle=\ttfamily\small,
	keywordstyle=\color{blue}\bfseries,
	commentstyle=\color{gray},
	stringstyle=\color{teal},
	numbers=left,
	numberstyle=\tiny\color{gray},
	stepnumber=1,
	numbersep=5pt,
	backgroundcolor=\color{lightgray!10},
	frame=single,
	breaklines=true,
	breakatwhitespace=false,
	tabsize=4,
	showspaces=false,
	showstringspaces=false,
	showtabs=false,
	captionpos=b,
	morekeywords={long long, vector, pair}
}

\newtcolorbox{mybox}[2][]{
	colback=#1!5!white,
	colframe=#1!75!black,
	fonttitle=\bfseries,
	title={#2}
}

\title{\Huge \textbf{二分查找算法完全指南}}
\author{算法学习笔记}
\date{\today}

\begin{document}
	
	\maketitle
	
	\tableofcontents
	\newpage
	
	\section{二分查找的核心概念}
	
	\subsection{什么是二分查找？}
	
	二分查找（Binary Search）是一种在有序序列中快速查找特定元素的算法。其核心思想是"减而治之"——每次都将搜索范围缩小一半。
	
	\begin{mybox}{blue}{形象理解}
		想象你在翻一本词典找某个单词：
		\begin{enumerate}
			\item 打开到中间位置
			\item 如果目标单词应该在前面，就把后半部分丢掉
			\item 如果目标单词应该在后面，就把前半部分丢掉
			\item 重复这个过程，直到找到目标
		\end{enumerate}
	\end{mybox}
	
	\subsection{为什么时间复杂度是O(log n)？}
	
	每次查找都能排除一半的数据：
	\begin{itemize}
		\item 第1次：剩 n/2 个数据
		\item 第2次：剩 n/4 个数据
		\item 第3次：剩 n/8 个数据
		\item ...
		\item 第k次：剩 n/2^k 个数据
	\end{itemize}
	
	当 $n/2^k = 1$ 时，$k = \log_2 n$，因此最多需要 $\log_2 n$ 次比较。
	
	\subsection{二分查找的关键——"二段性"}
	
	\textbf{重要理解：}二分查找的本质不是数据"有序"，而是数据具有"二段性"——我们可以用一个条件把数据分成"满足条件"和"不满足条件"的两段。
	
	\begin{mybox}{yellow}{二段性示例}
		\begin{itemize}
			\item 在有序数组中，以某个值为界，分为"小于目标值"和"大于等于目标值"
			\item 在旋转数组中，以旋转点为界，分为"较大的有序段"和"较小的有序段"
			\item 在峰值问题中，分为"在上升坡上"和"在下降坡上"
		\end{itemize}
	\end{mybox}
	
	\section{整数二分查找的两种经典模板}
	
	\subsection{模板一：查找第一个满足条件的元素（左边界）}
	
	\textbf{适用场景：}找到第一个满足条件的元素位置
	
	\begin{mybox}{green}{模板一核心记忆}
		\textbf{"左边界四步法":}
		\begin{enumerate}
			\item 初始化：$l = 0, r = n - 1$
			\item 循环条件：$while(l < r)$
			\item 中点计算：$mid = (l + r) / 2$（向下取整）
			\item 判断更新：
			\begin{itemize}
				\item 如果 $mid$ 满足条件：$r = mid$（不排除mid，向左收缩）
				\item 如果 $mid$ 不满足条件：$l = mid + 1$（排除mid，向右走）
			\end{itemize}
		\end{enumerate}
	\end{mybox}
	
	\begin{lstlisting}[caption={模板1：查找第一个 >= x 的元素}]
		#include <iostream>
		#include <vector>
		using namespace std;
		
		// 在有序数组arr中，查找第一个 >= x 的元素下标
		// 如果不存在，返回 -1
		int findFirstGreaterOrEqual(vector<int>& arr, int x) {
			int l = 0, r = arr.size() - 1;
			
			while (l < r) {
				int mid = l + (r - l) / 2;  // 等价于 (l + r) / 2，但防止溢出
				
				if (arr[mid] >= x) {
					// mid满足条件（>=x），可能是答案，往左找更小的
					r = mid;
				} else {
					// mid不满足条件（<x），mid肯定不是答案
					l = mid + 1;
				}
			}
			
			// while结束时 l == r，这个位置就是答案
			if (arr[l] >= x) return l;
			return -1;  // 所有元素都 < x
		}
		
		// 使用示例
		int main() {
			vector<int> arr = {1, 3, 3, 3, 5, 7, 9};
			int x = 3;
			int pos = findFirstGreaterOrEqual(arr, x);
			cout << "第一个 >= " << x << " 的位置: " << pos << endl;
			// 输出: 第一个 >= 3 的位置: 1
			return 0;
		}
	\end{lstlisting}
	
	\subsection{模板二：查找最后一个满足条件的元素（右边界）}
	
	\textbf{适用场景：}找到最后一个满足条件的元素位置
	
	\begin{mybox}{orange}{模板二核心记忆}
		\textbf{"右边界四步法":}
		\begin{enumerate}
			\item 初始化：$l = 0, r = n - 1$
			\item 循环条件：$while(l < r)$
			\item 中点计算：$mid = (l + r + 1) / 2$（向上取整，防止死循环）
			\item 判断更新：
			\begin{itemize}
				\item 如果 $mid$ 满足条件：$l = mid$（不排除mid，向右收缩）
				\item 如果 $mid$ 不满足条件：$r = mid - 1$（排除mid，向左走）
			\end{itemize}
		\end{enumerate}
	\end{mybox}
	
	\begin{lstlisting}[caption={模板2：查找最后一个 <= x 的元素}]
		#include <iostream>
		#include <vector>
		using namespace std;
		
		// 在有序数组arr中，查找最后一个 <= x 的元素下标
		// 如果不存在，返回 -1
		int findLastLessOrEqual(vector<int>& arr, int x) {
			int l = 0, r = arr.size() - 1;
			
			while (l < r) {
				int mid = l + (r - l + 1) / 2;  // 注意：+1是关键，防止死循环
				
				if (arr[mid] <= x) {
					// mid满足条件（<=x），可能是答案，往右找更大的
					l = mid;
				} else {
					// mid不满足条件（>x），mid肯定不是答案
					r = mid - 1;
				}
			}
			
			if (arr[l] <= x) return l;
			return -1;  // 所有元素都 > x
		}
		
		// 使用示例
		int main() {
			vector<int> arr = {1, 3, 3, 3, 5, 7, 9};
			int x = 3;
			int pos = findLastLessOrEqual(arr, x);
			cout << "最后一个 <= " << x << " 的位置: " << pos << endl;
			// 输出: 最后一个 <= 3 的位置: 3
			return 0;
		}
	\end{lstlisting}
	
	\subsection{为什么模板二需要 mid = (l + r + 1) / 2？}
	
	这是为了防止死循环。考虑 $l = 0, r = 1$ 的情况：
	
	\begin{mybox}{red}{死循环分析}
		\begin{itemize}
			\item \textbf{如果使用} $mid = (l + r) / 2$：
			\begin{itemize}
				\item $mid = (0 + 1) / 2 = 0$
				\item 如果满足条件，执行 $l = mid = 0$
				\item 区间还是 $[0, 1]$，没变化！进入死循环
			\end{itemize}
			\item \textbf{如果使用} $mid = (l + r + 1) / 2$：
			\begin{itemize}
				\item $mid = (0 + 1 + 1) / 2 = 1$
				\item 如果满足条件，执行 $l = mid = 1$
				\item 区间变成 $[1, 1]$，退出循环
			\end{itemize}
		\end{itemize}
	\end{mybox}
	
	\section{浮点数二分查找}
	
	浮点数二分不需要考虑整数边界问题，只需要控制精度。
	
	\begin{lstlisting}[caption={浮点数二分：求平方根}]
		#include <iostream>
		#include <cmath>
		using namespace std;
		
		// 计算x的平方根，精确到1e-6
		double sqrtBinary(double x) {
			double l = 0, r = max(1.0, x);  // 注意处理x<1的情况
			
			// 方法1：控制绝对误差
			while (r - l > 1e-8) {  // 精度要求
				double mid = (l + r) / 2;
				if (mid * mid >= x) {
					r = mid;
				} else {
					l = mid;
				}
			}
			return l;
		}
		
		// 方法2：固定循环次数（更常用）
		double sqrtBinary2(double x) {
			double l = 0, r = max(1.0, x);
			
			for (int i = 0; i < 100; i++) {  // 循环100次足够精确
				double mid = (l + r) / 2;
				if (mid * mid >= x) {
					r = mid;
				} else {
					l = mid;
				}
			}
			return l;
		}
		
		int main() {
			double x = 2.0;
			cout << "sqrt(2) = " << sqrtBinary(x) << endl;
			cout << "实际值: " << sqrt(2) << endl;
			// 输出: sqrt(2) ≈ 1.414214
			return 0;
		}
	\end{lstlisting}
	
	\section{STL中的二分查找函数}
	
	C++ STL提供了现成的二分查找函数：
	
	\begin{lstlisting}[caption={STL二分查找函数使用示例}]
		#include <iostream>
		#include <algorithm>
		#include <vector>
		using namespace std;
		
		int main() {
			vector<int> arr = {1, 3, 3, 3, 5, 7, 9};
			
			// 1. binary_search：判断元素是否存在
			bool exists = binary_search(arr.begin(), arr.end(), 3);
			cout << "3是否存在: " << (exists ? "是" : "否") << endl;
			
			// 2. lower_bound：查找第一个 >= x 的位置（左边界模板）
			auto left = lower_bound(arr.begin(), arr.end(), 3);
			cout << "第一个 >= 3 的位置: " << (left - arr.begin()) << endl;
			cout << "对应的值: " << *left << endl;
			
			// 3. upper_bound：查找第一个 > x 的位置
			auto right = upper_bound(arr.begin(), arr.end(), 3);
			cout << "第一个 > 3 的位置: " << (right - arr.begin()) << endl;
			cout << "对应的值: " << *right << endl;
			
			// 4. equal_range：同时获取左右边界
			auto range = equal_range(arr.begin(), arr.end(), 3);
			cout << "3的范围: [" << (range.first - arr.begin()) 
			<< ", " << (range.second - arr.begin()) << ")" << endl;
			
			// 总结：如果我们要找3出现的次数
			int count = right - left;
			cout << "3出现的次数: " << count << endl;
			
			return 0;
		}
	\end{lstlisting}
	
	\section{深入理解：二段性与check函数}
	
	\subsection{什么是好的check函数？}
	
	二分查找的核心是设计一个\textbf{check函数}，使得数据具有"二段性"：
	
	\begin{mybox}{blue}{check函数设计原则}
		\begin{itemize}
			\item \textbf{单调性：}前面一段不满足，后面一段满足（或反之）
			\item \textbf{明确性：}对于任意位置，能明确判断是否满足条件
			\item \textbf{目标性：}check函数的临界点就是我们想找的答案
		\end{itemize}
	\end{mybox}
	
	\subsection{如何构造check函数？}
	
	\begin{lstlisting}[caption={check函数的设计模式}]
		// 模式1：找第一个满足条件的（左边界）
		// check函数：判断是否满足目标性质
		bool check(vector<int>& arr, int mid, int target) {
			return arr[mid] >= target;  // 满足条件的性质
		}
		
		int binarySearchLeft(vector<int>& arr, int target) {
			int l = 0, r = arr.size() - 1;
			while (l < r) {
				int mid = l + (r - l) / 2;
				if (check(arr, mid, target)) {
					r = mid;  // 满足条件，往左找
				} else {
					l = mid + 1;  // 不满足，往右找
				}
			}
			return l;
		}
		
		// 模式2：找最后一个满足条件的（右边界）
		int binarySearchRight(vector<int>& arr, int target) {
			int l = 0, r = arr.size() - 1;
			while (l < r) {
				int mid = l + (r - l + 1) / 2;
				if (check(arr, mid, target)) {
					l = mid;  // 满足条件，往右找
				} else {
					r = mid - 1;  // 不满足，往左找
				}
			}
			return l;
		}
	\end{lstlisting}
	
	\section{经典例题详解}
	
	\subsection{例题1：在排序数组中查找元素的第一个和最后一个位置}
	
	\textbf{题目描述：}给定一个升序排列的整数数组和一个目标值，找出目标值在数组中的开始位置和结束位置。
	
	\begin{lstlisting}[caption={例题1：查找元素的起始和结束位置}]
		#include <iostream>
		#include <vector>
		using namespace std;
		
		vector<int> searchRange(vector<int>& nums, int target) {
			if (nums.empty()) return {-1, -1};
			
			// 第一次二分：找第一个 >= target 的位置（左边界）
			int l = 0, r = nums.size() - 1;
			while (l < r) {
				int mid = l + (r - l) / 2;
				if (nums[mid] >= target) {
					r = mid;
				} else {
					l = mid + 1;
				}
			}
			
			// 检查是否找到了target
			if (nums[l] != target) return {-1, -1};
			int left_bound = l;
			
			// 第二次二分：找最后一个 <= target 的位置（右边界）
			r = nums.size() - 1;  // l不变，只重置r
			while (l < r) {
				int mid = l + (r - l + 1) / 2;  // 注意+1
				if (nums[mid] <= target) {
					l = mid;
				} else {
					r = mid - 1;
				}
			}
			
			return {left_bound, l};
		}
		
		int main() {
			vector<int> nums = {5, 7, 7, 8, 8, 10};
			int target = 8;
			vector<int> result = searchRange(nums, target);
			
			cout << "数组: ";
			for (int num : nums) cout << num << " ";
			cout << endl;
			cout << "查找目标: " << target << endl;
			cout << "范围: [" << result[0] << ", " << result[1] << "]" << endl;
			// 输出: [3, 4]
			
			return 0;
		}
	\end{lstlisting}
	
	\subsection{例题2：寻找峰值}
	
	\textbf{题目描述：}峰值元素是指其值大于左右相邻值的元素。给定一个数组，返回任意一个峰值元素的索引。
	
	\textbf{解题思路：}虽然数组无序，但具有二段性：
	\begin{itemize}
		\item 如果 $nums[mid] < nums[mid+1]$，那么mid在上升坡上，峰值一定在右侧
		\item 否则，mid在下降坡或峰顶，峰值在左侧（包括mid）
	\end{itemize}
	
	\begin{lstlisting}[caption={例题2：寻找峰值}]
		#include <iostream>
		#include <vector>
		using namespace std;
		
		int findPeakElement(vector<int>& nums) {
			int l = 0, r = nums.size() - 1;
			
			while (l < r) {
				int mid = l + (r - l) / 2;
				
				// 比较mid和mid+1，判断在上升坡还是下降坡
				if (nums[mid] > nums[mid + 1]) {
					// mid在下降坡，峰值在左侧（包括mid）
					r = mid;
				} else {
					// mid在上升坡，峰值在右侧（不包括mid）
					l = mid + 1;
				}
			}
			
			return l;  // l就是峰值位置
		}
		
		int main() {
			vector<int> nums = {1, 2, 1, 3, 5, 6, 4};
			int peak = findPeakElement(nums);
			
			cout << "数组: ";
			for (int num : nums) cout << num << " ";
			cout << endl;
			cout << "峰值位置: " << peak << ", 峰值大小: " << nums[peak] << endl;
			// 可能的输出: 峰值位置: 5, 峰值大小: 6
			
			return 0;
		}
	\end{lstlisting}
	
	\subsection{例题3：搜索旋转排序数组}
	
	\textbf{题目描述：}一个原本升序的数组在某个点进行了旋转，搜索目标值的位置。
	
	例如：[4,5,6,7,0,1,2] 在原数组[0,1,2,4,5,6,7]的基础上旋转了4个位置。
	
	\textbf{解题思路：}二段性体现在——总能找到一段是有序的，判断target是否在这段中。
	
	\begin{lstlisting}[caption={例题3：搜索旋转排序数组}]
		#include <iostream>
		#include <vector>
		using namespace std;
		
		int searchRotated(vector<int>& nums, int target) {
			int l = 0, r = nums.size() - 1;
			
			while (l <= r) {  // 注意：这里是 <=
				int mid = l + (r - l) / 2;
				
				if (nums[mid] == target) {
					return mid;  // 找到了
				}
				
				// 判断哪边是有序的
				if (nums[l] <= nums[mid]) {
					// 左半部分有序
					if (nums[l] <= target && target < nums[mid]) {
						// target在左半部分
						r = mid - 1;
					} else {
						// target在右半部分
						l = mid + 1;
					}
				} else {
					// 右半部分有序
					if (nums[mid] < target && target <= nums[r]) {
						// target在右半部分
						l = mid + 1;
					} else {
						// target在左半部分
						r = mid - 1;
					}
				}
			}
			
			return -1;  // 没找到
		}
		
		int main() {
			vector<int> nums = {4, 5, 6, 7, 0, 1, 2};
			int target = 0;
			
			cout << "旋转数组: ";
			for (int num : nums) cout << num << " ";
			cout << endl;
			cout << "查找目标: " << target << endl;
			cout << "位置: " << searchRotated(nums, target) << endl;
			// 输出: 位置: 4
			
			return 0;
		}
	\end{lstlisting}
	
	\subsection{例题4：找第一个错误版本}
	
	\textbf{题目描述：}有n个版本[1, n]，已知在某个版本之后都是错误的。用最少的调用次数找到第一个错误的版本。
	
	\begin{lstlisting}[caption={例题4：第一个错误版本}]
		#include <iostream>
		using namespace std;
		
		// 模拟的API：判断版本是否正确
		bool isBadVersion(int version) {
			// 假设从版本4开始都是错误的
			return version >= 4;
		}
		
		int firstBadVersion(int n) {
			int l = 1, r = n;
			
			while (l < r) {
				int mid = l + (r - l) / 2;
				
				if (isBadVersion(mid)) {
					// 当前版本是错误的，第一个错误版本可能是mid或更早
					r = mid;
				} else {
					// 当前版本正确，第一个错误版本在后面
					l = mid + 1;
				}
			}
			
			return l;  // l就是第一个错误版本
		}
		
		int main() {
			int n = 10;
			int bad = firstBadVersion(n);
			
			cout << "总版本数: " << n << endl;
			cout << "第一个错误版本: " << bad << endl;
			// 输出: 第一个错误版本: 4
			
			return 0;
		}
	\end{lstlisting}
	
	\subsection{例题5：计算x的平方根}
	
	\textbf{题目描述：}实现int sqrt(int x)函数，计算并返回x的平方根，只保留整数部分。
	
	\begin{lstlisting}[caption={例题5：x的平方根（整数版本）}]
		#include <iostream>
		using namespace std;
		
		int mySqrt(int x) {
			if (x == 0) return 0;
			
			int l = 1, r = x;
			
			while (l < r) {
				int mid = l + (r - l + 1) / 2;  // 使用右边界模板
				
				// 判断mid的平方是否 <= x
				if (mid <= x / mid) {  // 用除法避免溢出
					// mid^2 <= x，满足条件，可能还有更大的答案
					l = mid;
				} else {
					// mid^2 > x，不满足，答案在左边
					r = mid - 1;
				}
			}
			
			return l;
		}
		
		int main() {
			int x1 = 8, x2 = 16, x3 = 2147395599;
			
			cout << "sqrt(" << x1 << ") = " << mySqrt(x1) << endl;  // 2
			cout << "sqrt(" << x2 << ") = " << mySqrt(x2) << endl;  // 4
			cout << "sqrt(" << x3 << ") = " << mySqrt(x3) << endl;  // 46339
			
			return 0;
		}
	\end{lstlisting}
	
	\subsection{例题6：在D天内完成所有任务}
	
	\textbf{题目描述：}传送带上有n个包裹，第i个包裹重量为weights[i]。要在D天内按顺序发送所有包裹，求传送带的最小运载能力。
	
	\textbf{解题思路：}二分答案——能力值越大，所需天数越少。找到第一个满足天数<=D的能力。这就是典型的"对答案进行二分"。
	
	\begin{lstlisting}[caption={例题6：传送带运载能力}]
		#include <iostream>
		#include <vector>
		#include <algorithm>
		using namespace std;
		
		// 计算以capacity的运载能力需要多少天
		int daysNeeded(vector<int>& weights, int capacity) {
			int days = 1;  // 至少需要1天
			int current_sum = 0;
			
			for (int weight : weights) {
				if (current_sum + weight > capacity) {
					// 当前包裹放不下了，需要新的一天
					days++;
					current_sum = weight;
				} else {
					current_sum += weight;
				}
			}
			
			return days;
		}
		
		int shipWithinDays(vector<int>& weights, int D) {
			// 运载能力的范围：至少是最大包裹重量，最大是所有包裹总重量
			int left = *max_element(weights.begin(), weights.end());
			int right = 0;
			for (int w : weights) right += w;
			
			// 二分查找最小运载能力（左边界模板）
			while (left < right) {
				int mid = left + (right - left) / 2;
				
				if (daysNeeded(weights, mid) <= D) {
					// 以mid的运力能在D天内完成，说明mid可能太大了
					// 尝试更小的运力
					right = mid;
				} else {
					// 以mid的运力不能在D天内完成，必须增大运力
					left = mid + 1;
				}
			}
			
			return left;
		}
		
		int main() {
			vector<int> weights = {1, 2, 3, 4, 5, 6, 7, 8, 9, 10};
			int D = 5;
			
			int capacity = shipWithinDays(weights, D);
			cout << "包裹重量: ";
			for (int w : weights) cout << w << " ";
			cout << endl;
			cout << "需要在 " << D << " 天内完成" << endl;
			cout << "最小运载能力: " << capacity << endl;
			// 输出: 15
			
			return 0;
		}
	\end{lstlisting}
	
	\section{二分答案——一种重要的解题思想}
	
	\subsection{什么时候使用二分答案？}
	
	\begin{mybox}{blue}{二分答案的典型场景}
		当一个问题满足以下条件时，可以尝试二分答案：
		\begin{enumerate}
			\item \textbf{答案具有单调性}：答案越大（或越小），越容易满足条件
			\item \textbf{可以快速验证}：对于给定的答案，能在较短时间内判断是否可行
			\item \textbf{求最值}：通常求"最大值的最小值"或"最小值最大值"
		\end{enumerate}
	\end{mybox}
	
	\subsection{二分答案的通用框架}
	
	\begin{lstlisting}[caption={二分答案的通用模板}]
		// 假设我们要找"满足条件的最小值"
		int binaryAnswer(int left, int right) {
			while (left < right) {
				int mid = left + (right - left) / 2;  // 左边界模板
				
				if (check(mid)) {
					// mid满足条件，尝试更小的值
					right = mid;
				} else {
					// mid不满足条件，需要更大的值
					left = mid + 1;
				}
			}
			return left;  // 最小的满足条件的答案
		}
		
		// check函数：判断某个答案是否可行
		bool check(int answer) {
			// 根据问题要求，判断以answer作为答案是否满足条件
			// 返回true表示满足，false表示不满足
		}
	\end{lstlisting}
	
	\section{总结：二分查找的核心要点}
	
	\begin{mybox}{green}{必须掌握的关键点}
		\begin{enumerate}
			\item \textbf{理解二段性}：二分查找的核心不是"有序"，而是能找到一种属性将数据分成两段
			\item \textbf{模板选择}：
			\begin{itemize}
				\item 找左边界（第一个满足条件的）：$mid = (l+r)/2, \; r = mid, \; l = mid+1$
				\item 找右边界（最后一个满足条件的）：$mid = (l+r+1)/2, \; l = mid, \; r = mid-1$
			\end{itemize}
			\item \textbf{防止死循环}：使用右边界模板时，$mid$计算必须$+1$
			\item \textbf{防止溢出}：使用 $mid = l + (r-l)/2$ 而不是 $(l+r)/2$
			\item \textbf{二分答案}：很多最优解问题可以通过二分答案转化为判定问题
		\end{enumerate}
	\end{mybox}
	
	\section{练习题推荐}
	
	\begin{enumerate}[leftmargin=*]
		\item \textbf{基础训练}：
		\begin{itemize}
			\item 查找插入位置（LeetCode 35）
			\item 找出有序数组中的单一元素（LeetCode 540）
			\item 寻找比目标字母大的最小字母（LeetCode 744）
		\end{itemize}
		\item \textbf{变形训练}：
		\begin{itemize}
			\item 搜索旋转排序数组 II（LeetCode 81）
			\item 寻找旋转排序数组中的最小值（LeetCode 153）
			\item 找出峰值（LeetCode 162）
		\end{itemize}
		\item \textbf{二分答案}：
		\begin{itemize}
			\item 吃香蕉的速度（LeetCode 875）
			\item 分割数组的最大值（LeetCode 410）
			\item 制作m天所需的花束（LeetCode 1482）
		\end{itemize}
	\end{enumerate}
	
\end{document}